\documentclass[11pt]{article}
\usepackage{template}
\usepackage{amsmath,amssymb,microtype,tikz}
\usetikzlibrary{arrows.meta,positioning}
\setlength{\emergencystretch}{2em}
\SetKeys{date=2026-10-26,vol=013,title={每日有趣新闻｜世界、科学与数学},author={ChatGPT 编辑部},note={\small 采编窗口：2026年10月24—26日（北京时间26日清晨截稿）。周末arXiv无常规新批次，数学选自23日上架的论文；注意区分选举夜估算与正式计票、政治人物主张与独立证实、预印本与期刊论文。}}
\renewcommand{\newsitem}[3]{\par\medskip{\large\bfseries\color{newsblue}#1\par}{\footnotesize\color{newsgray}#2\par}\smallskip\color{black}#3\par\medskip{\color{newsgray}\rule{\linewidth}{0.25pt}}\par}
\newcommand{\newslead}[3]{\par\medskip{\LARGE\bfseries\color{newsblue}#1\par}\vspace{1mm}{\small\color{newsgray}#2\par}\smallskip\color{black}#3\par\medskip{\color{newsgray}\rule{\linewidth}{0.45pt}}\par}
\newcommand{\newssub}[1]{\par\smallskip\noindent{\bfseries\color{newsblue}#1}\quad}
\begin{document}
\begin{documents}

\newslead{塞尔维亚选举之夜：学生运动如何走进议会，政权更替又意味着什么？}{国际政治头条｜10月25日议会选举；美联社现场报道；席位仍为选举夜估算}{
塞尔维亚10月25日举行议会选举，长期掌握国家权力的亚历山大·武契奇在投票结束后公开承认败选。美联社报道，武契奇在本党选举夜活动上表示，根据塞尔维亚前进党（SNS）的估算，由学生运动支持的“学生名单”预计取得250席议会中的135至140席，而前进党约获105席。这个数字来自参选方当晚宣布的估计，\textbf{不是中央选举机关已认证的最终结果}；各党席位、争议选区及组阁程序仍需后续核实。即便如此，败选承认本身已经改变政治局势：贝尔格莱德等城市出现大规模庆祝，人们开始讨论新多数如何把抗议诉求转化为可执行的制度改革。\cite{APSerbia}

\begin{center}
\begin{tikzpicture}[x=0.043cm,y=0.47cm]
\fill[newsblue!75] (0,1.6) rectangle (137.5,2.25);
\fill[newsgray!60] (0,0.3) rectangle (105,0.95);
\node[anchor=west,font=\scriptsize] at (0,2.63) {学生名单：预计135--140席（中点示意）};
\node[anchor=west,font=\scriptsize] at (0,1.33) {SNS：约105席};
\node[anchor=west,font=\scriptsize] at (0,-0.23) {议会共250席；数据为武契奇选举夜估算，非正式计票};
\draw[newsgray,dashed] (125,0) -- (125,2.4);
\node[anchor=south,rotate=90,font=\tiny] at (125,2.4) {125席过半线};
\end{tikzpicture}
\par{\scriptsize\color{newsgray}图1｜编辑部据美联社10月25日报道绘制。学生名单用135--140席中点137.5作示意，不是最终席位。}\end{center}

\newssub{从车站坍塌到选票。}2024年11月，诺维萨德一座火车站的混凝土雨棚坍塌，造成16人死亡。学生与其他抗议者认为，公共工程监管、承包和政治裙带关系需要受到调查，随后形成持续两年的大规模示威。政府支持者则强调稳定、发展和国家安全，武契奇把部分反对活动描述为外国势力介入。两类政治叙事并不等于事故调查结论：事故直接技术原因、招标与施工责任、是否存在腐败，必须依赖司法证据，而不能仅凭任何一方的竞选语言判断。\cite{APSerbia}

与传统反对党结盟不同，学生运动采取了较特别的组织方式。美联社指出，“学生名单”由大学生推动，但实际候选人包括专家、农民、运动员和艺术家，而不是由运动学生直接填满选票。这种安排试图解决一个老问题：抗议可以聚集对现状的不满，却未必自然产生有治理经验的候选人；如果把社会动员与候选人遴选分开，运动可能更容易吸纳不同背景的人。与此同时，名单内部的政策分歧、领导机制与问责方式，也会在选后显露出来。\cite{APSerbia}

\newssub{选举胜利不是制度改革的完成时。}武契奇此前担任总统，9月辞去总统职务，转而寻求总理职位。他和前进党过去14年间积累的行政网络、地方组织、媒体影响力及公共资源分配能力，不会因为一个选举夜自动消失。学生支持者希望重建司法独立、媒体自由与反腐制度；但要推动改革，还须依法完成计票、处理选举申诉、形成议会多数、组建内阁并通过预算和法律。选举监测人员已经报告了涉嫌买票、拍摄选票及对选民施压的情况，这些报告应得到独立调查，不能被胜选庆祝掩盖。\cite{APSerbia}

对外政策也有多重约束。塞尔维亚是欧盟候选国，同时与俄罗斯维持特殊关系，国内还存在科索沃地位问题与复杂的民族历史记忆。欧盟委员会主席冯德莱恩祝贺学生运动并表示愿意支持法治和媒体改革，但这不是欧盟入盟谈判立即完成的承诺。新政府即便有改革意愿，也需处理经济增长、能源、外交及地方政治的现实成本。外部支持、内部民意和行政执行能力是三件不同的事。\cite{APSerbia}

这场选举之所以值得在国际新闻中占据头条，不在于可以提前宣布塞尔维亚已走向何种政治终点，而在于它展示了社会抗议如何尝试进入代议政治。接下来最有解释力的指标将是正式计票结果、申诉裁决、组阁谈判、司法调查和具体改革法案，而不是任何阵营的胜利口号。
}

\newslead{鸭绿江新桥终于通车：朝中接触的外交象征与物流现实}{东北亚外交深读｜金正恩10月24日访问丹东；中朝官方发布与美联社交叉核对}{
10月24日，朝鲜领导人金正恩经新鸭绿江大桥进入中国辽宁丹东，随后与中共中央政治局常委王沪宁会面，并参访纪念场所和工厂。美联社依据新华社、朝中社以及韩国统一部的信息确认了访问和过桥行程。这是金正恩约一年以来首次访华，发生在朝中纪念朝鲜战争相关历史节点之际。两国官方强调传统友谊与务实合作，朝鲜媒体将此行置于历史叙事之中；报道这些表态时应说明其发布者，而不能把政治修辞当作双方已经签订的新经济协议。\cite{APNorthKorea,GuardianBridge}

\newssub{一座桥为何拖延十多年？}新鸭绿江大桥由中方投入大量资金建设，中国一侧较早完工，朝鲜一侧道路与通关设施长期未能投入使用，直到今年10月才正式开放。大桥设计上能够承载重型货车与集装箱运输，理论上改善跨境物流效率，但基础设施的物理通车与贸易量增长并不是一回事。运输能力还取决于边检通关、铁路公路衔接、货物许可、支付安排、国际制裁以及两国企业实际订单。没有这些数据，不能因为领导人乘车经过就推断双边贸易即将增加某个百分比。\cite{APNorthKorea,GuardianBridge}

\newssub{朝中俄关系并非简单三角形。}朝鲜近年与俄罗斯加强军事与经济合作，美联社报道朝鲜向俄方提供人员与常规武器以支持俄乌战争；中国则仍是朝鲜重要的贸易伙伴和外交沟通渠道。对平壤而言，增加与北京的交通联系可以降低过度依赖单一伙伴的风险；对北京而言，维持边境稳定、经济往来与朝鲜半岛局势可控具有现实意义。但双方利益并非在所有问题上完全一致，尤其涉及联合国制裁、核问题与地区安全安排时，更需要区分外交合作与政策一致。\cite{APNorthKorea}

韩国政府希望中国在朝美接触中发挥建设性作用。美国总统特朗普与金正恩是否会在11月的APEC会议期间会面，仍属外交猜测，不能从新桥通车推出峰会已经确定。历史纪念、交通工程和高层会晤在同一周发生，使事件具备象征性；然而其后续影响应通过跨境运输数据、公开协议和谈判进展观察，而不是把象征直接写成成果。\cite{APNorthKorea}

\newssub{照片的证据边界。}美联社页面展示新华社发布的过桥照片，也展示朝鲜官方提供的行程照片，并特别注明独立记者未获准进入部分现场。两种照片能够证明官方公开展示了何种场景，却不能独立核实所有行程细节。相关图片还涉及新闻摄影许可，本刊因此不直接转载，而采用文字与来源说明；本期其他文章使用原创信息图。
}

\newsitem{铁矿石进口增加、港口库存更快上升：中国与澳大利亚矿商在争什么？}{产业经济｜路透社10月26日专栏，数据截至10月23日；专栏分析与原始统计分开}{
如果只看海关数字，中国今年前九个月铁矿石进口量达9.5497亿吨，同比增长4.1\%，似乎说明钢铁需求依然强劲。但路透社10月26日刊发的商品市场专栏指出，钢铁生产与房地产建设仍显疲弱，进口增加的一部分进入了港口库存：SteelHome跟踪的10月23日当周港口库存约1.5223亿吨，比上年同期增加14\%。同期中国粗钢产量同比下降2.9\%。进口、产量、库存并不矛盾，它们分别描述进入市场的原料、当期工业加工与尚未被消耗的存量。\cite{ReutersOre}

\begin{center}
\begin{tikzpicture}[x=0.052cm,y=0.45cm]
\fill[newsblue!65] (0,1.6) rectangle (104.1,2.2);
\fill[newsgray!65] (0,0.4) rectangle (114,1);
\node[anchor=west,font=\scriptsize] at (0,2.55) {进口量指数：104.1（同比+4.1\%）};
\node[anchor=west,font=\scriptsize] at (0,1.35) {港口库存指数：114（同比+14\%）};
\draw[densely dashed] (100,0.1)--(100,2.3);
\node[anchor=north,font=\tiny] at (100,0.1) {上年同期=100};
\end{tikzpicture}
\par{\scriptsize\color{newsgray}图2｜两项同比指数，不是同一绝对量。进口为2026年前九个月累计；港口库存为10月23日当周。来源：海关数据及SteelHome，经路透社整理。}\end{center}

\newssub{买方集中与卖方集中同时存在。}中国约占全球海运铁矿石购买量的四分之三，但澳大利亚大型矿商提供中国约三分之二的铁矿石需求。采购方有巨大的市场规模，矿商则掌握难以在短期替代的高品位矿源、物流网络与稳定交货能力。中国矿产资源集团（CMRG）尝试集中采购与协调议价，以争取更有利的定价指数和合同条款。路透社援引知情人士称，CMRG近期与力拓达成供应协议，仍与Fortescue存在谈判分歧；具体合同价格没有公开，因此不能宣称中方已经获得确定的折扣幅度。\cite{ReutersOre}

价格变化也不能全部归功于谈判。新加坡交易所相关铁矿石合约10月23日约每吨90.60美元，比5月11日的年内高点111.91美元低19\%；钢铁需求下降、库存积累、供给预期和市场情绪都可能参与定价。几内亚西芒杜项目未来或能增加替代供应，但矿山、铁路、港口与产能爬坡需要时间。把长期供应链重组写成短期价格胜负，容易高估任何一方的谈判能力。\cite{ReutersOre}

这份报道还提醒读者区分新闻与评论：上述进口、库存、价格是可核对的统计数据；关于集中采购取得多少谈判优势的判断，则属于路透社专栏作者的市场分析，并非所有矿商共同确认的结论。未来值得关注的，是公开合同条件、钢厂利润与库存周转速度，而不只是单日铁矿石报价。
}

\newsitem{乌克兰称美方弹道导弹预警恢复：战场信息依赖如何进入外交谈判？}{俄乌与美国外交｜路透社10月25日报道；关键事实为泽连斯基公开陈述}{
乌克兰总统泽连斯基10月25日接受电视采访时表示，美国关于俄罗斯弹道导弹袭击的预警信息已经完全恢复。他还称，卫星通信服务Starshield此前受到限制，近日已恢复到可支持高速移动条件下使用的状态；他把两项变化与自己在奥斯陆会见美国特使威特科夫联系起来。路透社报道了上述说法，但这不等于美国政府、服务运营商已分别公开证实每一项技术细节。尤其是报道中的“每小时420公里”描述的是使用条件下的移动速度，\textbf{不是网络传输速率}。\cite{ReutersUkraine}

预警信息对防空系统的价值在于争取反应时间。发现发射、估算弹道、分配拦截资源和发布民防警报，是连续而不同的环节；外部情报若突然减少，防御方可能面临更大的不确定性。通信服务又关系到部队之间的数据传输和行动协调。两者都是现代战争中难以仅靠单一国家独立提供的能力，因此援助政策与战场行动会相互牵连。\cite{ReutersUkraine}

华盛顿要求基辅停止袭击俄罗斯炼油设施，理由之一是担心能源价格；泽连斯基则提出俄方停止袭击乌克兰电力设施、美国提供额外防空能力或放宽相关技术使用条件等交换选项。这些是谈判立场，不是已经达成的停火条款。泽连斯基还表示，不接受在11月技术谈判中讨论限制乌克兰军队规模或强制中立，并坚持领土问题须由国家领导人层级讨论。能否举行峰会、双方是否接受监督机制，目前都没有确定答案。\cite{ReutersUkraine}
}

\newsitem{量子信息的化学新候选：室温稳定的三价锆分子为何值得关注？}{化学与量子物理｜\emph{Nature Chemistry} 10月23日同行评审论文}{
量子计算研究中的“量子比特”并不一定要由超导电路或囚禁离子实现。某些顺磁性分子拥有未成对电子，其自旋可以形成可操纵的量子态。难题在于：自旋会与周围原子核、分子振动和环境磁噪声耦合，导致相干性消失。对于过渡金属配合物，设计稳定的配位环境并控制轨道成分，是同时解决化学稳定性与自旋可控性的途径。\cite{ZrQudits}

10月23日发表于《自然·化学》的研究报告了可在室温下稳定分离的三价锆芳氧基配合物。锆是相对丰富的金属，但Zr(III)配合物往往容易二聚或发生歧化反应。研究者利用体积较大的芳氧基配体构造近似三角平面的配位环境，使单占分子轨道具有明显的锆$4d_{z^2}$与$5s$混合特征。理论计算估计该轨道的$s$成分约25--27\%，与较强的各向同性超精细耦合相联系。\cite{ZrQudits}

实验通过脉冲电子顺磁共振观察自旋弛豫与相干操控：在5 K条件下，部分样品自旋--晶格弛豫时间$T_1$达到约87毫秒，去相位相关时间最长约123微秒，并观察到电子及电子--核复合跃迁的相干操控。这里必须区分三个层次：\textbf{室温下化学上稳定}，不意味着\textbf{室温下具有同样长的量子相干时间}，更不意味着已经造出可扩展的通用量子计算机。论文的贡献是扩展了可供研究的分子量子态平台。\cite{ZrQudits}

对于配位化学而言，研究还展示了几何构型与电子结构之间的关系：并非只能靠稀有元素才能得到可操纵的量子态，也可以利用配体空间位阻和轨道混成重新安排电子的局部环境。今后需要比较不同温度下的相干时间、器件集成、单分子寻址和多量子态耦合，才能判断这种材料在实际量子信息硬件中的位置。
}

\newsitem{数学拾趣｜把1到537分成六组，每组都不能出现$x+y=z$}{组合数论｜arXiv:2610.31362，10月22日提交、23日上架；含可检查证书，尚非同行评审定论}{
先从一道小题开始：把整数$1,2,3,4$分成两组，要求每一组都不存在三个数$x,y,z$满足$x+y=z$，其中允许$x=y$。一种可行分法是$\{1,4\}$与$\{2,3\}$；前一组没有$1+1=2$，因为2不在该组，后一组的$2+2=4$也不在该组。这样不含内部加法等式的集合称为\textbf{无和集}（sum-free set）。\cite{Schur}

\begin{center}
\begin{tikzpicture}[scale=0.65]
\foreach \n/\c in {1/newsblue,2/newsgray,3/newsgray,4/newsblue} {
  \pgfmathtruncatemacro{\xx}{\n-1}
  \fill[\c!35] (\xx*1.15,0) rectangle ++(0.96,0.85);
  \draw[newsgray] (\xx*1.15,0) rectangle ++(0.96,0.85);
  \node at (\xx*1.15+0.48,0.43) {\n};
}
\node[anchor=west,font=\scriptsize] at (0,-0.5) {蓝组\{1,4\}；灰组\{2,3\}；两组内都没有$x+y=z$。};
\end{tikzpicture}
\par{\scriptsize\color{newsgray}图3｜编辑部原创最小示例：$\{1,2,3,4\}$的两色无和划分。它只解释定义，不是537的证明证书。}\end{center}

对给定颜色数$k$，Schur数$S(k)$是能把$\{1,\dots,n\}$分成$k$个无和集的最大$n$。问题不是把整数分成“大小差不多”的组，而是同时避开所有同色的加法三元组。随着$n$增大，可选的整数越来越多，限制也越来越密集。对于六种颜色，长期使用的下界是536；陈志扬、朱立耕和龙凡在10月22日提交预印本，给出$1$至$537$的一个明确六色划分，从而主张$S(6)\ge537$。作者还报告找到131个在所采用等价关系下彼此不同的划分，并提供完整划分、检查脚本与证书文件。\cite{Schur}

\newssub{“下界”与“精确值”差别很大。}要证明$S(6)\ge537$，只需要给出一个有效划分并逐项验证它没有违例；这是一种\textbf{存在性证明}。要证明$S(6)=537$，还必须证明538个整数无论怎样分成六组都会出现违例，这是另一种\textbf{全称性结论}，通常困难得多。新预印本没有因为把下界推进一格，就自动解决精确值。附带的机器可读证书使前半个主张具有可复核性，但本刊没有独立运行所有检查脚本，因此只能准确报道作者给出的构造和验证途径。\cite{Schur}

为什么增加一个整数也值得报道？组合极值问题的进展不一定体现在数量级的跃迁。构造一个此前未知的有效划分，可以揭示隐藏的排列规律、对称性与搜索策略，并成为更大规模计算的起点。论文指出，已有递推构造还可据此更新若干更高阶Schur数的下界；这些数字仍应理解为由构造推出的下界，而非已确定的精确答案。\cite{Schur}
}

\newsitem{另一则数学进展：康威的“每两条边恰好相交一次”猜想更接近上界}{图论与拓扑｜arXiv:2610.31589，10月22日提交；尚未独立审查}{
在平面上画一个图，允许边弯曲，但要求任意两条边\textbf{恰好相遇一次}：它们可以在公共端点相接，也可以在内部交叉，却不能既共享端点又再相交一次。这样的图画法称为thrackle。康威猜想，若图有$n$个顶点，那么边数最多$n$条。直觉上“让每一对边都相遇”似乎有利于增加交叉，实际上它同时施加了极强的几何限制。\cite{Thrackle}

Dániel Gábor Simon提交的七页预印本声称，对于$n\ge4$，可证明边数不超过$\lfloor4(n-1)/3\rfloor$。这一上界仍大于康威猜想要求的$n$，因此\textbf{并未证明原猜想}，但它进一步收紧了允许的边数。论文将此前对至多12个顶点图的计算机辅助核验，与平面和射影平面嵌入上的放电法结合。放电法的核心思想是给局部结构分配可守恒的“电荷”，再利用局部不等式排除不可能的全局配置。\cite{Thrackle}

这份工作目前只是arXiv预印本。本刊没有重做其计算机枚举或拓扑论证，故把结论写为作者提出的新证明主张；后续应观察专家核查、代码复现和正式发表情况。
}

\newsitem{以色列的“加密货币绕过制裁”设想遭央行与经济学家反对}{国际经济与选举｜路透社10月25日报道；选举在本刊截稿时尚未举行}{
以色列总理内塔尼亚胡在电视采访中谈及若遭受更广泛国际制裁，可能探索加密货币和区块链技术作为经济应对手段。此前英国、法国和加拿大宣布限制进口来自约旦河西岸以色列定居点的商品。以色列银行（央行）公开反驳这一思路，认为加密货币不能成为绕开制裁的可靠路径，主张优先维护与主要贸易伙伴的关系。超过一百名经济学家也发表公开信，警告金融隔离可能冲击投资、就业和民生。\cite{ReutersCrypto}

加密货币网络可以在技术层面转移数字资产，却不能凭空提供石油、设备、保险、国际清算银行或合法的跨境物流。经济制裁涉及的不仅是支付通道，还包括贸易对手、银行合规、货物运输与金融市场准入。把“某种资产可以点对点转移”直接等同于“一个国家可以绕过全面经济制裁”，忽略了实体经济的约束。另一方面，经济学家的警告是对潜在政策后果的预测，并非制裁已经发生后的统计结果。\cite{ReutersCrypto}

报道发生在以色列议会选举前夕，执政党与反对派围绕国家安全和经济管理相互批评。选举尚未举行，不能将民调推测写成选举结果。本刊关注的是提议本身的金融可行性，以及政策争论中“技术方案”如何被用于回答更广泛的外交问题。
}

\newsitem{Rust能防数据竞争，却不保证驱动程序的操作顺序正确}{软件工程研究简报｜arXiv:2610.27325，10月19日提交；为避免遗漏补收，非当日新论文}{
Rust的所有权与借用检查能在安全代码中阻止许多内存错误和数据竞争，但\textbf{逻辑竞态}仍可能存在：两个线程分别持有合法锁、各自操作也合法，组合后的执行顺序却破坏业务不变量。例如先检查某设备状态、再执行依赖该状态的操作，期间另一线程可能改变设备状态。\cite{XRust}

Xu等人提出XRust，通过分析操作之间的逻辑先后关系，自动构造并发测试，比较并发执行是否等价于两种串行顺序之一。作者称在Linux 7.3-rc2的十个Rust驱动中构建1654个测试用例，发现两处此前未知的逻辑竞态，其中一处已获维护者确认并有修复进入相应子系统树。这是作者预印本中的实验报告，不能扩展成“所有Rust驱动都有同类漏洞”，也不能误读为Rust内存安全保证失效。\cite{XRust}
}

\end{documents}
\vspace{1mm}
{\small\bfseries\color{newsblue}参考资料与核查说明}\par
{\footnotesize 采编截止为北京时间2026年10月26日上午。数学论文均为未完成本刊独立复核的预印本；政治报道中的选举夜数字、政府声明与分析意见均作了性质区分。三张配图由编辑部使用TikZ原创，非现场照片。}\par
\begingroup\footnotesize
\renewcommand{\refname}{参考文献}\bibliographystyle{plainurl}
\bibliography{refs}
\endgroup
\end{document}